\section{Existing Range Filters}
\label{sec:existing-work}

We survey the five practical range filters most directly comparable
to our work from the past seven years.
We group them by the technique
that determines their space--accuracy trade-off, contrast each with
the Goswami lower bound (Section~\ref{sec:bounds}), and use the
grouping to position our contributions.

\paragraph{Trie-based: SuRF.}
The Succinct Range Filter~\citep{zhang2018surf} stores the set of keys
as a Fast Succinct Trie --- a level-order unary degree sequence (LOUDS)
encoding with dense and sparse layers, $10$ bits per trie node,
$4$--$15\times$ faster than earlier succinct tries on point-query
workloads. SuRF answers a range query by walking
the trie down to the boundary of the query interval, and in a RocksDB
deployment accelerates closed-seek queries by up to $5\times$ over the
Bloom-filter baseline. The trie skeleton gives no clean parameterisation
by $\varepsilon$: SuRF-Base costs about $10$ bits per key for $64$-bit
integers and $14$ bits per key for email addresses regardless of the
target FPR, and the authors themselves note that ``query keys are
almost always correlated to the stored keys'', driving the empirical
FPR up to $4\%$ on integer datasets and $25\%$ on
emails~\citep[\S3.1]{zhang2018surf}. SuRF does not position itself
against an information-theoretic range-filter space bound and does not
cite the Goswami construction.

\paragraph{Bloom cascade: Rosetta.}
Rosetta~\citep{luo2020rosetta} stores one Bloom filter per key-prefix
length, so a range query of length $R$ decomposes into at most
$2 \log_2 R$ dyadic intervals, each initiating a Bloom probe at the
root of its sub-tree. The cascade's space is within a factor of $1.44$
of the Goswami lower bound under uniform per-level
allocation~\citep[\S3.1]{luo2020rosetta} and is reported $2$--$5\times$
faster than SuRF and up to $40\times$ faster than the default RocksDB
on a $50$M-key benchmark at $22$ bits per key. The authors are explicit
about the cost: ``due to the multiple probes required (for every
prefix) Rosetta has a high probe cost''~\citep[\S1]{luo2020rosetta},
mitigated by allocating the per-level bit budget from runtime query
statistics gathered between LSM compactions.

\paragraph{Learned: SNARF.}
SNARF~\citep{vaidya2022snarf} fits a linear-spline model to the
empirical key CDF and stores a Golomb-coded sparse bit array (with an
Elias--Fano alternative) over the predicted positions. The published
FPR is approximately $1/K$ where $K$ scales with the bit budget; on
the configurations covered by their evaluation SNARF reaches up to
$50\times$ better FPR than SuRF and $100\times$ better FPR than
Rosetta at matched memory. The model accuracy, however, is
distribution-dependent: on the OSM dataset, whose sorted keys contain
large empty contiguous ranges and abrupt jumps, SNARF's FPR--space
curve degrades visibly and reaches $10^{-4}$ only at substantially
higher memory than on other real-world
datasets~\citep[\S6.1.2, Fig.~5(B)]{vaidya2022snarf}. SNARF additionally has no
worst-case bound under correlated queries: ``when a query end point
is consistently close to an actual key but the query does not include
a key, it may yield a false positive in SNARF and
SuRF''~\citep[\S4.1]{vaidya2022snarf}. In our SOSD~Facebook benchmark
with workload-correlated queries (Chapter~\ref{chap:evaluation}) the
empirical FPR floors around $10^{-3}$ regardless of additional memory.

\paragraph{Practical realisation of Goswami: Grafite.}
Grafite~\citep{grafite2024} is the closest related work to this thesis.
It is explicitly framed as a practical, simplified instantiation of the
Goswami construction: a locality-preserving hash
$h(x) = (q(\lfloor x/r \rfloor) + x) \bmod r$ with pairwise-independent
$q$ maps the universe down to a smaller range, the resulting hash
codes are stored using Elias--Fano coding, and a range query reduces
to a single predecessor probe. By the authors' own framing, Grafite is
``the only range filter to date to achieve a robust and predictable
false positive rate across all combinations of datasets, query
workloads, and range sizes''~\citep[\S1]{grafite2024}. Its space
matches the Goswami lower bound to an additive $2n$ bits, and on
uniform data it is reported $6.7$--$10.3\times$ faster to construct
than Rosetta. What Grafite does not do --- and what we add in this
thesis --- is exploit the structure of the input keys: its
space--accuracy trade-off is by design workload-agnostic. Our
truncation hash (contribution~3) replaces Grafite's locality-preserving
rotation, and Scan-ARE (contribution~4) layers distribution-aware
segmentation on top of the same ERE-style backend.

\paragraph{Dynamic extension: Memento.}
Memento~\citep{eslami2024memento} is the contemporaneous SIGMOD~2024
filter that adds dynamicity --- insertions, deletions, and growth ---
to the static state of the art. It is built on top of a Rank-and-Select
Quotient Filter: each universe partition stores a fingerprint together
with a contiguous list of fixed-length key suffixes (``mementos''),
all packed via Robin Hood hashing. Memento reports false-positive
behaviour and per-query cache-miss count comparable to Grafite at
matched memory. The authors explicitly note that static range filters
suffice for the immutable files of an LSM-Tree
deployment~\citep[\S1]{eslami2024memento}; Memento targets workloads
where the key set itself is mutable --- single-global-filter
LSM-trees, B-Tree indexes, and systems that mix transactional and
analytical workloads. Because our thesis stays
within the static setting, Memento is orthogonal rather than a
head-to-head competitor.

\paragraph{Positioning.}
The space of practical range filters splits along a single axis: does
the FPR guarantee hold under any query workload, or only under some
distribution assumption? Trie-based and learned designs fall on the
heuristic side, but for different reasons. SuRF has no
parameterisation by $\varepsilon$ at all: its space is fixed by the
trie skeleton, and its FPR is whatever falls out of the suffix budget,
empirically up to $25\%$ on correlated workloads. SNARF is
parameterised by $\varepsilon$ but loses the guarantee under correlated
queries, where the linear-spline model and the surrounding key context
inflate the FPR. The robust side is occupied by Rosetta (within a
factor of $1.44$ of the lower bound, $O(\log R)$ Bloom probes per
query) and Grafite (within an additive $2n$ bits, single predecessor
probe). Memento extends the robust side to dynamic workloads at
parity on the static axis. None of the five exploits the
cluster-and-gap structure of real LSM keys; this regularity is left
on the table. The thesis fills the gap: we keep the worst-case
guarantee of the Grafite line and add a distribution-aware mechanism
that beats the bound on realistic inputs.

\paragraph{Contributions of this thesis.}
We start from Goswami's two-layer construction --- locality-preserving
hash followed by an ERE backend --- and improve it along three
orthogonal axes, producing a family of filters that match the Goswami
lower bound on adversarial inputs and beat it on realistic ones.
\begin{enumerate}
    \item \textbf{Single-vector ERE backend (Chapter~\ref{chap:ere}).}
    The two metadata vectors $D_1, D_2$ of Goswami's ERE are replaced
    by a single succinct vector $D$, reducing metadata size by
    $24\%$ --- a saving predicted from the Poisson bucket-occupancy
    model and confirmed empirically.

    \item \textbf{Adaptive bucket search (Chapter~\ref{chap:ere}).}
    The theoretically $O(1)$ Weak Prefix Search index inside each ERE
    bucket is replaced by adaptive binary/linear search on packed
    suffixes, removing $59$~bits of auxiliary metadata per key and
    reducing per-query latency by $13$--$30\times$ at LSM bucket
    sizes.

    \item \textbf{Truncation hash (Chapter~\ref{chap:hash}).}
    For sparse inputs we replace Goswami's locality-preserving
    rotation with the truncation map $h(x) = \lfloor x/2^t \rfloor$,
    which contracts the per-key collision zone from
    $\mathcal{L} \cdot 2^t$ to $\mathcal{L} + 2^t$ and saves
    $\log_2 \mathcal{L}$ bits per key with $O(1)$ query cost.

    \item \textbf{Distribution-aware segmentation: Scan-ARE
    (Chapter~\ref{chap:hybrid}).}
    A single-pass 1D-DBSCAN partitions the input into dense clusters
    and sparse tails. Dense clusters are handed to an ERE filter
    (exact, zero false positives); sparse points fall back to the
    truncation hash. On clustered SOSD data Scan-ARE drives the FPR
    below $10^{-8}$ at $\sim\!11$ bits per key. At matched memory the
    strongest prior art either degrades under workload correlation
    (SNARF, floor $\sim\!10^{-3}$ in our measurements) or, in the
    distribution-free line (Grafite, Rosetta), only reaches an FPR
    of $\sim\!10^{-2}$.
\end{enumerate}

\paragraph{Roadmap.}
Chapter~\ref{chap:ere} develops the single-vector ERE and the adaptive
bucket-search backend (contributions~1--2). Chapter~\ref{chap:distributions}
characterises the key distributions we target, preparing the ground for
contributions~3--4. Chapter~\ref{chap:hash} formalises and compares
the SODA and Truncation hash designs (contribution~3).
Chapter~\ref{chap:hybrid} introduces the Scan-ARE segmentation
(contribution~4). Chapter~\ref{chap:evaluation} reports the empirical
comparison against SuRF, SNARF, Rosetta, and Grafite on SOSD and
synthetic workloads. Chapter~\ref{chap:conclusion} concludes.
